\section{Scaling Laws: de la intuición empírica a una convicción central}

\subsection{Origen: la revelación de 2014-2017}

Dario se incorporó a Baidu en 2014, bajo la mentoría de Andrew Ng, para trabajar en reconocimiento de voz. En ese entonces, la opinión dominante en el mundo académico era que «nos falta el algoritmo correcto», pero Dario, como recién llegado, planteó una pregunta sencilla: ¿qué pasaría si el modelo se hiciera más grande, se agregaran más datos y se entrenara durante más tiempo?

\begin{knowledgebox}{El origen de la Scaling Hypothesis}
Dario no es el único creyente en el Scaling. Ilya Sutskever tiene una frase célebre y casi zen: \textit{«The thing you need to understand about these models is they just want to learn.»}

El Bitter Lesson de Rich Sutton y la Scaling Hypothesis de Gwern apuntaban en la misma dirección. Pero el verdadero punto de inflexión fueron los resultados de GPT-1 en 2017, que convencieron a Dario de que el lenguaje era el mejor vehículo para el Scaling.
\end{knowledgebox}

\begin{importantbox}{Analogía de la reacción química}
El Scaling requiere ampliar linealmente y al mismo tiempo tres «reactivos»: una red más grande, más datos y más cómputo. \textit{«Si solo amplías uno de ellos sin ampliar los demás, la reacción se detiene.»} Esto explica por qué muchos investigadores, en las primeras etapas, no vieron el efecto del Scaling: solo ampliaron una o dos de esas dimensiones.
\end{importantbox}

\subsection{¿Por qué funciona el Scaling? El ruido 1/f y la distribución de cola larga}

Dario recurre a su formación en física (en el posgrado estudió biofísica) para explicar el mecanismo subyacente del Scaling. Introduce un concepto clave: el \textbf{ruido 1/f}, es decir, la idea de que los procesos naturales producen distribuciones de cola larga.

El lenguaje natural tiene una estructura jerárquica que va de lo simple a lo complejo:
\begin{enumerate}
    \item Distribución de frecuencia de palabras (los patrones más comunes)
    \item Reglas gramaticales básicas
    \item Estructura a nivel de oración
    \item Coherencia temática a nivel de párrafo
    \item Ideas novedosas y razonamiento complejo
\end{enumerate}

\textit{«Primero capturan correlaciones muy simples… y después viene una cola muy larga con otros patrones más complejos.»} La suavidad de la distribución de cola larga se refleja directamente en el descenso suave de la curva de loss: por eso las Scaling Laws muestran una relación de ley de potencia tan regular.

Idea clave: este marco no solo explica \textit{por qué} funciona el Scaling, sino también \textit{por qué el progreso parece gradual y no repentino}: porque el modelo va «excavando» progresivamente a lo largo de la distribución de cola larga hacia patrones cada vez más raros.

\subsection{¿Dónde está el techo de la inteligencia?}

Dario distingue tres niveles de techo:

\textbf{Por debajo del nivel humano}: no existe techo. \textit{«Los humanos somos capaces de entender estos patrones»}, así que en teoría el modelo también debería poder alcanzarlos.

\textbf{Más allá del nivel humano: depende del campo}. En biología puede haber un margen enorme: \textit{«los humanos apenas logramos entender la complejidad de la biología»}, un campo con disciplinas muy fragmentadas donde es difícil integrar una visión de conjunto. En cambio, en campos como la ciencia de materiales o el conflicto humano, el techo podría estar más cerca del nivel humano.

\textbf{El techo institucional}: esta es la opinión más singular de Dario:

\begin{warningbox}{El verdadero cuello de botella no es la inteligencia, sino las instituciones}
El sistema de ensayos clínicos es un ejemplo excelente: incluye tanto burocracia innecesaria (procesos excesivamente conservadores) como mecanismos de protección razonables (la seguridad de los pacientes). Aunque la IA lograra diseñar un fármaco perfecto, este igual tendría que pasar por ese proceso lento.

\textit{«Somos demasiado lentos y demasiado conservadores.»} Dario considera que las instituciones humanas (los sistemas burocráticos, los sistemas regulatorios, la inercia organizacional) podrían ser un cuello de botella mayor que las limitaciones puramente de inteligencia.
\end{warningbox}

\subsection{Posibles obstáculos para el Scaling}

Dario analiza uno por uno los factores que podrían impedir que el Scaling continúe:

\textbf{Limitaciones de datos}: los datos disponibles en internet son limitados y su calidad está bajando (contenido repetido, spam de SEO, contenido generado por IA). Pero hay dos salidas:
\begin{itemize}
    \item \textbf{Datos sintéticos}: usar el modelo para generar sus propios datos de entrenamiento. AlphaGo Zero es el ejemplo definitivo: puro autojuego, cero datos humanos, y aun así alcanzó un nivel sobrehumano.
    \item \textbf{Modelos de razonamiento}: el Chain-of-Thought combinado con aprendizaje por refuerzo es, en esencia, también una forma de generar datos sintéticos.
\end{itemize}

\textbf{El costo del cómputo}: entrenar los modelos de frontera actuales cuesta cerca de \$1,000 millones; el próximo año se llegará a decenas de miles de millones de dólares; en 2026 se superarán los \$10,000 millones; y la ambición para 2027 es un clúster de cómputo del orden de \$100,000 millones. \textit{«Creo que todo esto va a suceder. La determinación para construir esta capacidad de cómputo es enorme.»}

\textbf{Limitaciones arquitectónicas}: podría hacer falta un método de optimización o una arquitectura completamente nuevos. Pero, por ahora, no hay indicios de que la arquitectura Transformer haya llegado a una limitación fundamental.

\textbf{Que los modelos dejen de progresar}: esto es lo más difícil de predecir; tal vez exista alguna razón desconocida que detenga las mejoras.

La evidencia más contundente en contra: \textbf{el avance de SWE-bench}, que pasó de 3\% en enero de 2024 a 50\% en octubre, en tan solo 10 meses. \textit{«Si seguimos extrapolando… en pocos años estos modelos superarán a los humanos con el nivel profesional más alto.»}

\begin{importantbox}{El «efecto película» del Scaling}
\textit{«Ya he visto esta película suficientes veces… en cada etapa del Scaling siempre han existido argumentos en contra, y cada vez han sido superados.»} Dario reconoce que las Scaling Laws son una regularidad empírica y no una ley física, pero elige «apostar a que va a continuar».
\end{importantbox}

\subsection{Resumen de la sección}

El mecanismo central de las Scaling Laws es la propiedad de cola larga del lenguaje: mientras más grande es el modelo, más patrones raros puede capturar. Los tres «reactivos» (la red, los datos y el cómputo) deben ampliarse de manera sincronizada. Por ahora no hay ningún factor de obstáculo convincente, y el rápido avance de SWE-bench es la evidencia empírica más sólida. Pero el verdadero cuello de botella tal vez no esté en la inteligencia misma, sino en la velocidad de adaptación de las instituciones humanas.
